\documentclass[10pt,twocolumn,letterpaper]{article}
\usepackage[final]{cvpr}
\usepackage[T1]{fontenc}
\usepackage{times,amsmath,booktabs}
\usepackage{flushend}
\usepackage[hidelinks]{hyperref}
\begin{document}
\title{PitchClipers: Context-Aware Football Highlight Construction\newline with Player--Ball Evidence}
\author{Zaid Rafi}
\maketitle
\begin{abstract}
Football highlights must preserve an event's build-up and outcome while limiting viewing time. PitchClipers combines pretrained action spotting, asymmetric temporal decoding and player--ball evidence. We formulate context-based evaluation and implement threshold calibration. A controlled decoder experiment removes 32.7\% duplicate viewing time through interval merging, motivating separate measurement of action confidence and contextual sufficiency.
\end{abstract}

\section{Motivation and Related Work}
A goal timestamp omits the attacking build-up; a foul frame may omit the resulting decision. A useful highlight should preserve both event context and outcome with minimal unrelated footage. We ask: \emph{how can action proposals become context-complete highlights under a fixed viewing budget?}

SoccerNet-v2~\cite{sn} benchmarks timestamp-based spotting. CALF~\cite{calf} learns temporal context; NetVLAD++~\cite{netvlad} separates past and future features. ActionFormer~\cite{actionformer} predicts action boundaries through multiscale temporal attention. These methods localize actions, whereas highlight boundaries must also preserve explanatory context. Our contribution is a reproducible interval decoder and context-evaluation formulation, supported by spatial evidence from detection and tracking~\cite{yolo,bytetrack}.

\section{Computer-Vision Method}
\textbf{Temporal proposals and decoding.} Video frames sampled at 2 FPS are represented by frozen ResNet-152 descriptors, reduced to 512 dimensions using the released PCA transform, and processed by pretrained CALF for 17 action classes. Each proposal is $(t_i,c_i,s_i)$: timestamp, class and confidence. A confidence-filtered proposal generates
\begin{equation}
 I_i=[\max(0,t_i-b),\min(T,t_i+a)],
\end{equation}
where $T$ is video duration. The default $b=8$ and $a=5$ seconds retain asymmetric context. Overlapping intervals are merged to avoid repeating footage. Users inspect the candidates before export; the fixed-window decoder provides a reproducible baseline.

\textbf{Spatial evidence.} Football Analytics combines player/ball detections with ByteTrack~\cite{bytetrack} to provide trajectories, ball visibility and proximity-control cues. Missing ball observations remain explicit. The image-coordinate offside line supports geometric inspection, without establishing a rules-based offside decision.

\textbf{Calibration and learned extension.} The implemented calibrator supports class-threshold optimization by coordinate grid search over validation context F1. A proposed learnable selector, $q_\theta(I)=\sigma(\theta^\top z_I)$, would combine class, confidence, duration and tracking features, trained with regularized binary cross-entropy on human sufficiency labels. Match-separated splits would prevent overlapping-window leakage. Selector training remains future work.

\section{Context-Based Evaluation}
Annotators specify the shortest interval $g_j$ preserving each target event's build-up and outcome. Let $\mathcal C$ be the merged output clips, $G=\bigcup_jg_j$ and $U=\bigcup_{I\in\mathcal C}I$. Define
\begin{equation}
 R=\frac1N\sum_j\max_{I\in\mathcal C}\frac{|I\cap g_j|}{|g_j|},
 \qquad P=\frac{|U\cap G|}{|U|}.
\end{equation}
Coverage $R$ measures retained context; precision $P$ penalizes unrelated footage. The evaluator reports their harmonic mean, complete-context rate and redundancy. We additionally propose Context Utility $H=S\,2PR/(P+R)$, where $S$ is the fraction of clips containing a correctly labelled reference event; zero denominators yield zero. This evaluation applies automatically to any spotter given shared context annotations, which timestamp labels alone cannot supply.

\section{Decoder Sensitivity and Discussion}
With saved CALF proposals fixed, we varied confidence threshold $\tau$ and context $(b,a)$, isolating decoder effects. Table~\ref{tab:decoder} reports reel duration and overlap removed relative to concatenating unmerged windows.
\begin{center}
\small
\begin{tabular}{rrrrr}
\toprule
$\tau$ & $(b,a)$ [s] & Proposals & Reel [s] & Removed [s] \\
\midrule
0.05 & (4,3) & 2 & 11.5 & 2.5 \\
0.05 & (8,5) & 2 & 17.5 & 8.5 \\
0.20 & (8,5) & 1 & 13.0 & 0.0 \\
0.05 & (12,8) & 2 & 21.5 & 15.5 \\
\bottomrule
\end{tabular}
\refstepcounter{table}\label{tab:decoder}
\par\smallskip Table \thetable: Local decoder sensitivity on saved event proposals.
\end{center}
At fixed confidence, padding increases duration from 11.5 to 21.5 seconds without adding proposals. Default merging reduces 26.0 seconds to 17.5 seconds: 32.7\% less repeated viewing, preserving the temporal union. Raising the threshold removes one proposal and 4.5 seconds. Thus padding, filtering and merging separately affect available context, accepted evidence and viewing cost.

Additional footage can contain either useful build-up or unrelated play; duration alone cannot measure highlight quality. These local sensitivity results establish decoder behaviour. Evaluating contextual sufficiency requires held-out interval annotations and equal viewing budgets.

\textbf{Conclusion.} The interactive demo exposes event timelines, tracking overlays and MP4 reels. The next experiment compares fixed decoding and learned selection, with and without tracking features, using context scores at equal viewing duration.

{\small\begin{thebibliography}{6}
\bibitem{sn} A. Deli\`ege et al. SoccerNet-v2: A Dataset and Benchmarks for Holistic Understanding of Broadcast Soccer Videos. CVPRW, 2021.
\bibitem{calf} A. Cioppa et al. A Context-Aware Loss Function for Action Spotting in Soccer Videos. CVPR, 2020.
\bibitem{bytetrack} Y. Zhang et al. ByteTrack: Multi-Object Tracking by Associating Every Detection Box. ECCV, 2022.
\bibitem{yolo} J. Redmon et al. You Only Look Once: Unified, Real-Time Object Detection. CVPR, 2016.
\bibitem{netvlad} S. Giancola and B. Ghanem. \href{https://openaccess.thecvf.com/content/CVPR2021W/CVSports/papers/Giancola_Temporally-Aware_Feature_Pooling_for_Action_Spotting_in_Soccer_Broadcasts_CVPRW_2021_paper.pdf}{Temporally-Aware Feature Pooling for Action Spotting in Soccer Broadcasts}. CVPRW, 2021.
\bibitem{actionformer} C.-L. Zhang, J. Wu, and Y. Li. \href{https://www.ecva.net/papers/eccv_2022/papers_ECCV/html/7278_ECCV_2022_paper.php}{ActionFormer: Localizing Moments of Actions with Transformers}. ECCV, 2022.
\end{thebibliography}}
\end{document}
